\chapter{AMMs, liquidity provision and impermanent loss}\label{ch:amm}

\noindent The automated market maker replaced the order book with a function. Instead of
matching buyers and sellers, a pool holds reserves of two assets and quotes prices off a
curve. Understanding the curve---its price rule, its depth, and who bears what risk---is the
entry ticket to on-chain market design.

\section{The constant-product rule}\label{sec:cpmm}

A constant-product pool holds reserves $(x,y)$ and enforces $x\cdot y=k$ through every trade.
Its spot price is $p=y/x$; buying $x$ with $y$ pushes the state down the curve and raises $p$.
The quoted execution price is set by the curve's slope, not by any counterparty.

\begin{figure}[htbp]\centering
\begin{tikzpicture}[font=\scriptsize]
  \draw[->,slate] (0,0)--(6.4,0) node[below left=-1pt]{reserve $x$};
  \draw[->,slate] (0,0)--(0,4.6) node[above left=-1pt]{reserve $y$};
  \draw[accent,very thick] plot[domain=1.02:6,samples=60] (\x,{6/\x});
  \fill[ink] (2,3) circle (1.6pt) node[above right]{state $(x,y)$, price $y/x$};
  \fill[amber] (3,2) circle (1.6pt) node[below right]{after buy of $\Delta x$};
  \draw[->,ink] (2.15,2.75) -- (2.85,2.25);
  \node[text=slate,align=center] at (4.9,3.4) {$x\cdot y=k$\\ slope $\approx$ price\\ curvature $=$ depth};
\end{tikzpicture}

\caption{The constant-product curve: trades move the state along $x\cdot y=k$; the local
slope is the price, the curvature is the depth.}
\label{fig:cpmm}
\end{figure}

\begin{principle}[price is reserve ratio]
For CPMM, marginal price $=y/x$ exactly. A trade of $\Delta x$ gets
$\Delta y = \frac{\Delta x\,(1-f)\,y}{x+\Delta x(1-f)}$ with fee $f$---executed price worsens
with size because you are walking along the curve, which is why pool depth \emph{is}
liquidity.
\end{principle}

\begin{example}\label{ex:amm}
\textbf{Swaps and impact.} \texttt{amm\_swap.py} runs a pool of $x{=}1000$ tokens against
$y{=}50{,}000$ USD (price 50) at fee 0.3\%:
\begin{lstlisting}[numbers=none]
pool x=1000 y=50000 price=50.00
in    1: out     49.8  eff price  49.80  impact +0.40%  new px   49.90
in   10: out    493.6  eff price  49.36  impact +1.28%  new px   49.02
in   50: out  2374.1  eff price  47.48  impact +5.03%  new px   45.36
in  200: out  8312.5  eff price  41.56  impact +16.88%  new px   34.74
in  500: out  16633.3  eff price  33.27  impact +33.47%  new px   22.24
\end{lstlisting}
\begin{solution}
Impact grows roughly linearly for small trades (1 token $\to$ 0.4\%, 50 $\to$ 5\%) then
convexly: the last line buys at an average 33 vs spot 50. Note the double consequence---the
buyer got a bad fill \emph{and} moved the pool's spot to 22.24, which is now an arbitrage
invitation the moment any market still trades near 50.
\end{solution}
\end{example}

\begin{figure}[htbp]\centering
\begin{tikzpicture}
\begin{axis}[width=0.72\linewidth,height=0.36\linewidth,
    xlabel={trade size $\Delta x$ (of 1000 in pool)},ylabel={price impact (\%)},
    xmin=0,xmax=600,ymin=0,ymax=36,
    tick style={draw=none},axis line style={draw=slate},font=\scriptsize]
  \addplot[accent,very thick] file{figures/data/PI.tex};
  \node[font=\scriptsize,text=slate,align=left] at (axis cs:330,8)
    {convex: impact $\approx(\Delta x/x)$ for small trades};
\end{axis}
\end{tikzpicture}

\caption{Execution-price impact vs trade size in the same pool: convexity is the cost of
consuming depth.}
\label{fig:impact}
\end{figure}

\section{Arbitrage is the oracle}\label{sec:amm-arb}

Nothing inside the pool knows the true price. External markets do, and arbitrageurs buy the
cheap leg and sell the dear one until the pool's spot equals the market's. This has two
important readings: (i) pool prices track markets because bots force them to; (ii) \emph{the
LPer is always the other side} of that rebalancing trade. The LP sells what is rising and
buys what is falling---which is exactly why LP returns lag holding.

\begin{figure}[htbp]\centering
\begin{tikzpicture}[font=\scriptsize,every node/.style={align=center}]
  \node[draw=slate,fill=white,rounded corners=1pt,minimum width=3cm] (pool) at (0,0)
    {pool price\\ $y/x=45$};
  \node[draw=slate,fill=white,rounded corners=1pt,minimum width=3cm] (mkt) at (4.5,0)
    {market price\\ $50$};
  \node[draw=good,fill=goodbg,rounded corners=1pt,minimum width=3cm] (arb) at (2.25,1.7)
    {bot buys from pool,\\ sells on market};
  \draw[<->,ink,thick] (pool) -- (mkt);
  \draw[->,ink] (pool) -- (arb);
  \draw[->,ink] (mkt) -- (arb);
  \node[text=slate] at (2.25,-0.9) {trade continues until pool price $\approx$ market};
\end{tikzpicture}

\caption{The arbitrage channel: gap between pool spot and market price is closed by bots;
LPs fund the gap.}
\label{fig:ambarb}
\end{figure}

\begin{example}\label{ex:arbpnl}
\textbf{Arb PnL.} The pool of Example~\ref{ex:amm} sits at spot 22.24 after the oversized
trade while every external venue still quotes 50. A bot buys tokens from the pool until the
pool's marginal price reaches 50 (ignoring fees for the sketch). How many tokens close the
gap, and what does the rebalancing look like from the LP's seat?
\begin{solution}
Pool spot $=y'/x'=22.24$ with $x'y'\approx k=5\times10^{7}$: restoring spot 50 needs
$y/x=50$, i.e.\ reserves back to $(1000,\,50{,}000)$. The bot therefore \emph{buys}
$\approx500$ tokens out of the pool for $\approx16{,}667$ USD---an average fill of
$\approx33$ against a market price of 50. The LP delivered those tokens $17$/token below
market: about $\$8{,}300$ of value handed to the bot, roughly $16\%$ of pool TVL. That
transfer is the IL of fig.~\ref{fig:ilcurve} with a name on it.
\end{solution}
\end{example}

\section{Impermanent loss}\label{sec:il}

If the price ratio moves from 1 to $r=p_1/p_0$, the LP's pool value is
$V_{\mathrm{LP}}=2\sqrt{k}\sqrt{p_1}$-equivalent while the HODL value is
$V_{\mathrm{H}}=y_0+x_0p_1$. Their ratio depends only on $r$:
\begin{equation}\label{eq:il}
\frac{V_{\mathrm{LP}}}{V_{\mathrm{H}}}=\frac{2\sqrt{r}}{1+r},
\qquad \mathrm{IL}=\frac{2\sqrt{r}}{1+r}-1\le 0.
\end{equation}
IL is symmetric (a 2x move up hurts exactly as much as a 2x move down), scale-free (only the
\emph{ratio} matters), and path-independent at the endpoint (though each fee accrues on the
path).

\begin{figure}[htbp]\centering
\begin{tikzpicture}
\begin{axis}[width=0.72\linewidth,height=0.36\linewidth,
    xlabel={price ratio $r=p_1/p_0$},ylabel={LP vs HODL (\%)},
    xmin=0.15,xmax=4,ymin=-22,ymax=2,
    tick style={draw=none},axis line style=draw,axis lines=middle,
    xtick={0.25,0.5,1,1.5,2,3,4},font=\scriptsize]
  \addplot[accent,very thick] file{figures/data/I_IL.tex};
\end{axis}
\end{tikzpicture}

\caption{IL(\% ) vs price ratio: the $-5.7\%$ at $r=2$ and $r=0.5$---movement in either
direction costs the LP relative to holding.}
\label{fig:ilcurve}
\end{figure}

\begin{example}\label{ex:il}
\textbf{The IL table.} \texttt{il\_calc.py} evaluates \eqref{eq:il}:
\begin{lstlisting}[numbers=none]
r=0.25: IL=-20.00%
r=0.50: IL= -5.72%
r=0.80: IL= -0.62%
r=1.00: IL= +0.00%
r=1.25: IL= -0.62%
r=1.50: IL= -2.02%
r=2.00: IL= -5.72%
r=4.00: IL=-20.00%
\end{lstlisting}
\begin{solution}
The LP breaks even only if fees out-earn IL. At a 0.3\% fee and a $r=1.5$ move
(IL=$-2.0\%$), the pool needs $\approx0.224\times$ its TVL in daily volume \emph{every day
for a month} to compensate---most pools never see that, which is why LPing is a bet that
volume arrives faster than price moves.
\end{solution}
\end{example}

\begin{example}\label{ex:lp}
\textbf{LP vs HODL across seeds.} \texttt{lp\_sim.py} runs 40 GBM paths
($\sigma_{\mathrm{daily}}\approx3\%$, 180 days) with fee income at 8\% of TVL/day:
\begin{lstlisting}[numbers=none]
40 seeds: LP+fees-vs-hodl mean +2.48%  no-fees -1.93%
best +4.41%  worst -12.27%
\end{lstlisting}
\begin{solution}
Without fees the LP loses to holding by $\approx2\%$ on average---that's the IL from the
simulation's volatility. Adding a modest fee stream flips the mean positive but leaves a
$-12\%$ worst case: LPing is short convexity---you collect a stream of small fees against
occasional large directional givebacks. The honest description is \emph{selling volatility
to the market} and getting paid rent for it.
\end{solution}
\end{example}

\begin{figure}[htbp]\centering
\begin{tikzpicture}
\begin{axis}[width=0.86\linewidth,height=0.32\linewidth,
    xlabel={day},ylabel={LP minus HODL (\%)},
    xmin=0,xmax=180,ymin=-6,ymax=1,
    tick style={draw=none},axis line style={draw=slate},
    xtick={0,45,90,135,180},font=\scriptsize]
  \addplot[accent,thick] file{figures/data/LP_PATH.tex};
  \addplot[slate,dotted] coordinates {(0,0)(180,0)};
\end{axis}
\end{tikzpicture}

\caption{One seed's LP-minus-HODL path: the bleed is continuous whenever price trends in
either direction---the fee stream must outrun this drift.}
\label{fig:lppath}
\end{figure}

\section{Concentrated liquidity}\label{sec:cl}

Uniswap v3 lets an LP bind its liquidity to a price range $[a,b]$: inside the range the
position behaves like a CPMM with \emph{virtual} reserves scaled up by the leverage
$1/(1-\sqrt{a/b})$; outside it holds only one asset and earns nothing. \texttt{cl\_range.py}:
\begin{lstlisting}[numbers=none]
range +/-50%: capital multiplier ~2.4x
range +/-20%: capital multiplier ~5.4x
range +/-10%: capital multiplier ~10.5x
range +/-5%: capital multiplier ~20.5x
range +/-2%: capital multiplier ~50.5x
\end{lstlisting}
The multiplier multiplies both the fee rate \emph{and} the IL exposure per unit of
capital---a $\pm20\%$ range concentrates capital $\approx5\times$, so its IL is
$\approx5\times$ the full-range number in percentage-of-capital terms.

\begin{figure}[htbp]\centering
\begin{tikzpicture}[font=\scriptsize]
  \draw[->,slate] (0,0)--(9,0) node[below left=-1pt]{price};
  \draw[->,slate] (0,0)--(0,3.6);
  \draw[accent,very thick] plot[domain=0.4:8.6,samples=80] (\x,{1.8+1.6/(\x/4+0.4)});
  \draw[slate,dashed] (3,0)--(3,3.2); \draw[slate,dashed] (5,0)--(5,3.2);
  \draw[fill=accentsoft,draw=none] (3,0) rectangle (5,3.2);
  \draw[amber,very thick] plot[domain=3:5,samples=40] (\x,{1.8+1.6/(\x/4+0.4)});
  \node[text=slate] at (4,-0.5) {range $[a,b]$: all capital works here};
  \node[text=slate,align=left] at (7.4,2.6) {outside range:\\0 fees,\\all one asset};
\end{tikzpicture}

\caption{Concentrated liquidity: capital works only inside $[a,b]$; leaving the range stops
fees and leaves the position entirely in the underperforming asset.}
\label{fig:cl}
\end{figure}

\begin{example}\label{ex:range}
\textbf{Choosing a range.} ETH trades at 4000. You expect a quiet week: range $[3800,4200]$
($\pm5\%$, $\approx20\times$ capital efficiency). A competitor chooses $[3200,4800]$
($\pm20\%$, $\approx5\times$).
\begin{solution}
If price stays inside, you earn $\approx4\times$ the fees per unit of capital. On a move to
4300 you are out of range holding only the sold asset: no fees, full exposure to the move you
missed, and the decision to re-range costs gas and a new IL clock. Tight ranges are a
volatility call: you are compensated only while realized vol stays inside your band.
\end{solution}
\end{example}

\begin{checkpoint}[AMM sanity]
\begin{itemize}
\item Show \eqref{eq:il} from $V_{\mathrm{LP}}=2\sqrt{kp_1}$ and $V_{\mathrm{H}}=y_0+x_0p_1$.
\item From \texttt{amm\_swap.py}: at what trade size does effective price drop below the
post-trade spot? (Never---the curve guarantees execution price $\le$ new spot.)
\item Why does an LP never beat HODL on a one-way move? (IL is strictly negative for
$r\neq1$.)
\end{itemize}
\end{checkpoint}

\section{Recap}\label{sec:ammrecap}

The CPMM is a function that sells liquidity: $p=y/x$ is the price, curvature is the depth,
arbitrage is the oracle. LPing converts that function's rebalancing into a PnL: fee income
minus the IL of being on the wrong side of every move. Concentrated liquidity is the same
trade with a leverage knob on range width.

\begin{table}[htbp]\centering\small
\caption{Chapter toolbox.}
\begin{tabular}{@{}lll@{}}
\toprule
Concept & Formula/tool & Gotcha\\
\midrule
Spot price & $p=y/x$ & post-trade price, not mid\\
Execution & $\Delta y=\Delta x(1-f)y/(x+\Delta x(1-f))$ & convex impact in size\\
IL & $2\sqrt r/(1+r)-1$ & symmetric; ratio-only\\
CL multiplier & $1/(1-\sqrt{a/b})$ & scales fees \emph{and} IL\\
Breakeven & fee$\times$volume $>$ IL & volume must persist\\
\bottomrule
\end{tabular}
\end{table}

\begin{example}[The arb that restores the price]\label{ex:arb}
\textbf{Strategy.} External venue prices the token at 55 while the
pool sits at 50. Compute the exact trade that moves the pool to 55,
its cost, and the residual profit --- the ``free money'' behind
every AMM price.

\textbf{Solution.} On $x=1000,y=50000$ ($k=5\cdot10^7$),
\texttt{arb\_restore.py} prints:
\begin{verbatim}
buy 46.5374 token for 2440.44 quote
arb revenue 2559.56 cost 2440.44 profit 119.12
new spot 55.0000 (restored to external)
with 0.3% fee: cost 2447.79 profit 111.77 (fee shaves 7.34)
\end{verbatim}
The restoring trade buys $46.54$ tokens at an average of $52.44$ ---
cheaper than the 55 they resell for --- netting $\$119$ gross,
$\$111.8$ after the 0.3\% fee.

\textbf{Significance.} The AMM's price \emph{is} this arb: the pool
moves to fair value only because someone is paid to move it, and the
fee is skimmed off their profit, not yours. Notice the arb spent
$\$2447$ to extract $\$112$ --- a $4.6\%$ return on the flow, which
is why on thin pools the restoring arb is small and the misprice
can sit for blocks.
\end{example}

\begin{example}[Uptime times concentration]\label{ex:cluptime}
\textbf{Strategy.} A concentrated-liquidity position earns only
while price stays in range --- narrower ranges pay more per
in-range day but sit idle more often. Simulate one GBM path and
score three widths on \emph{fee-days equivalent}
(in-range days $\times$ concentration multiple).

\textbf{Solution.} \texttt{cl\_stress.py} prints:
\begin{verbatim}
same GBM path (sigma 2%/day, 300 days), three range widths:
range +/- 5%: in-range 188/300 (63%)  conc-mult 2.2x  fee-days-equiv 420
range +/-10%: in-range 252/300 (84%)  conc-mult 1.6x  fee-days-equiv 398
range +/-25%: in-range 292/300 (97%)  conc-mult 1.0x  fee-days-equiv 292
narrower width = fewer inside days x higher multiple
\end{verbatim}
The $\pm5\%$ range wins on fee-days ($420$ vs $292$) despite
being idle $37\%$ of the time --- the $2.2\times$ concentration
more than compensates at this volatility.

\textbf{Significance.} Range selection is not a mood: it is an
optimisation of uptime $\times$ multiplier, and the answer moves
with $\sigma$. The $\pm5\%$ position that wins here would lose to
a wider range in a choppier regime --- re-run the sim before you
re-run the range.
\end{example}

\section{Common mistakes}
\begin{itemize}
\item \textbf{``IL is only on downside.''} It is symmetric: $r=2$ and $r=0.5$ both cost
  $-5.7\%$. Any move, either direction, hurts vs holding.
\item \textbf{Forgetting the fee must outrun IL \emph{per unit time}.} A 0.3\% fee on slow
  volume will not cover a 5\% price move.
\item \textbf{Comparing LP return to zero.} The benchmark is HODL---being flat is not the
  alternative; holding the two assets is.
\item \textbf{Reading the multiplier as free leverage.} CL scales the IL surface, not just
  the fees: a tight range magnifies both.
\item \textbf{Assuming pool price = truth.} It is arbitrage'd to truth only while arbitrage
  is profitable after gas and latency---in thin moments it lags.
\item \textbf{Ignoring that impact is convex.} The same pool that gives 0.4\% on a small
  trade gives 33\% on a large one; quote impact, not fee.
\item \textbf{Thinking the pool ``discovers'' price.}
Nothing in the contract watches the outside world: an arbitrageur
moves the pool to fair value for a profit
(Example~\ref{ex:arb}). Your IL is their fee for the service.
\item \textbf{Ignoring the fee in arb math.} The 0.3\% skim
($\$7.34$ on the example) is small per trade but decides whether a
marginal misprice gets closed at all --- it sets the dead band
around fair value.\end{itemize}

\section{Exercises}
\begin{exercise}[easy]
A pool holds $(x,y)=(2000,100{,}000)$. Spot price and $k$?
\end{exercise}
\begin{solution}
$p=y/x=50$, $k=2\times10^{8}$.
\end{solution}

\begin{exercise}[easy]
Compute IL for $r=3$. How does it compare to $r=1/3$?
\end{exercise}
\begin{solution}
$2\sqrt3/4-1=-13.4\%$; symmetric: identical at $r=1/3$.
\end{solution}

\begin{exercise}[easy]
Buy 20 tokens of the $(1000,50{,}000)$ pool at fee 0.3\%. Effective price and new spot?
\end{exercise}
\begin{solution}
$\Delta y=19.94\times50000/1019.94=977.5$; eff $48.9$; new spot $49,022.5/1020=48.06$.
\end{solution}

\begin{exercise}[easy]
Why is a 2x price move equally bad for an LP as a halving?
\end{exercise}
\begin{solution}
IL depends on $\sqrt r + 1/\sqrt r$ symmetrically: the pool ends up holding more of the
cheaper asset in both cases.
\end{solution}

\begin{exercise}[easy]
What is the CL capital multiplier for a $\pm1\%$ range?
\end{exercise}
\begin{solution}
$1/(1-\sqrt{0.99/1.01})\approx1/(1-0.990)\approx100\times$---and so is the IL magnification.
\end{solution}

\begin{exercise}
Derive the marginal price of CPMM from $y=k/x$: show $-\partial y/\partial x = y/x$.
\end{exercise}
\begin{solution}
$y=kx^{-1}\Rightarrow dy/dx=-k/x^{2}=-y/x$; the executed price is the slope's magnitude.
\end{solution}

\begin{exercise}
A pool earns 0.3\% on daily volume equal to 15\% of TVL. Price drifts $r=1.3$ over a month.
Does the LP beat HODL?
\end{exercise}
\begin{solution}
Fees $\approx0.003\times0.15\times30=1.35\%$; IL at $r=1.3$: $2\sqrt{1.3}/2.3-1=-0.854\%$;
net $+0.5\%$---yes, marginally.
\end{solution}

\begin{exercise}
Show that the CPMM swap path in Example~\ref{ex:amm} keeps $x'y' \approx k$ up to the fee.
\end{exercise}
\begin{solution}
$(x+\Delta x)(y-\Delta y)=k$ holds for the fee-inclusive $\Delta x$ portion; the fee slice
accrues to reserves, raising $k$ slightly---the source of LP fee income.
\end{solution}

\begin{exercise}
An LP re-ranges a CL position weekly. Write the PnL decomposition that makes the strategy
comparable to full-range LP.
\end{exercise}
\begin{solution}
PnL = (fees while in range) $-$ (IL realized on exits) $-$ (gas per re-range) $\pm$ (asset
drift while out of range); all four terms scale with the multiplier, so the comparison is
not just fee rate.
\end{solution}

\begin{exercise}
Why does a large trade \emph{itself} create an arbitrage opportunity, and who pays for
closing it?
\end{exercise}
\begin{solution}
The trade moves the pool spot away from external prices; a bot buys the cheap leg on the
pool and sells on the venue---the LP pays, in the form of the rebalanced reserve mix.
\end{solution}

\begin{exercise}
Volume/TVL of a real pool is 8\%/day at 0.3\% fee. What annualized fee APR does that imply
before IL?
\end{exercise}
\begin{solution}
$0.08\times0.003\times365=8.8\%$ APR---which is why high turnover, not high headline APR,
is the LP metric that matters.
\end{solution}

\begin{exercise}
Sketch the LP-vs-HODL curve's asymptotes as $r\to\infty$: what does the LP end up holding?
\end{exercise}
\begin{solution}
IL $\to -100\%$: the pool is all $x$, the asset that went to zero value---the LP sold all
the appreciating asset along the way.
\end{solution}

\begin{exercise}[hard]
Prove that in a mean-reverting market with no fees, LP still loses to HODL for any nonzero
excursion---i.e.\ IL is not just a trending-market phenomenon. (Hint: IL is path-dependent
in value terms.)
\end{exercise}
\begin{solution}
For any path returning to $r=1$ the endpoint IL is zero---so in a \emph{perfectly} mean-
reverting path the loss is fees-only on net, but every round trip still transferred value to
the arbitrageurs who rebalanced the pool mid-path; the endpoint formula hides the
pathwise wealth transfer.
\end{solution}

\begin{exercise}[hard]
Model a $\pm10\%$ CL position as a barrier strategy. What does the position look like after
price exits the range---which option did you effectively write?
\end{exercise}
\begin{solution}
You hold 100\% of the depreciated asset---like being knocked into a covered position: the
CL LP is short a straddle with the range as strike band; exiting the band is the payoff of
having sold it.
\end{solution}

\begin{exercise}[hard]
Extend \texttt{lp\_sim.py}: measure the volume/TVL needed so that median (not mean)
LP-vs-HODL is positive. Report the number.
\end{exercise}
\begin{solution}
Sweep \texttt{fee\_daily\_turnover} across seeds; the median crossing is the answer
(typical runs: $\sim$4--6\% daily at these vols).
\end{solution}

\begin{exercise}[easy]
External price jumps to 55 while the pool sits at 50. Who trades
first: the LP trying to exit, or the arb? Why?
\end{exercise}
\begin{solution}
The arb --- the LP's exit does not move the pool's price (it
withdraws both sides proportionally), while the arb's buy
\emph{is} the price-restoring mechanism and pays a guaranteed
profit.
\end{solution}
\begin{exercise}[medium]
On the example pool, show the arb's average entry ($52.44$) sits
between the old spot ($50$) and the new ($55$). Why must it?
\end{exercise}
\begin{solution}
Each infinitesimal buy moves the spot up the curve; the average
price paid is the integral from 50 to 55 along it --- strictly
inside the interval, hence between the endpoints.
\end{solution}
\begin{exercise}[hard]
After the 0.3\% fee the arb nets $\$111.77$ on $\$2447.79$.
Compute the dead band: the external-price gap below which the arb
loses money, assuming the pool's spot is 50.
\end{exercise}
\begin{solution}
Profit $= m(x_0-x_1) - (y_1-y_0)/(1-f)$ vanishes when the external
price barely exceeds spot: near zero-gap the fee dominates, so the
band is $\approx$ one fee-tier around spot ($\sim$0.3\%--0.6\%).
Inside it, misprices persist unexploited.
\end{solution}

\section*{Notes}
AMMs made liquidity provision a mechanical, auditable game---and turned ``providing
liquidity'' into the most mis-sold product in DeFi. The math is one line ($x\cdot y=k$) and
one inequality (fees vs IL); everything else is frictions. For production rules (fee tiers,
oracle integration, tick math), the venue docs are canonical and change frequently---the
Uniswap docs cited below are the reference.

\section*{Further reading}
\begin{itemize}
\item Uniswap v2 docs, ``How Uniswap works'' --- the CPMM spec:

\url{https://docs.uniswap.org/contracts/v2/concepts/protocol-overview/how-uniswap-works}
\item Uniswap docs, ``Concentrated liquidity'' --- v3 range math:

\url{https://docs.uniswap.org/concepts/protocol/concentrated-liquidity}
\item Wikipedia, ``Automated market maker'' --- family taxonomy beyond CPMM:

\url{https://en.wikipedia.org/wiki/Automated_market_maker}
\item Wikipedia, ``Uniswap'' --- history and protocol structure:

\url{https://en.wikipedia.org/wiki/Uniswap}
\item DailyDeFi IL calculator --- sanity-check \eqref{eq:il} interactively:

\url{https://dailydefi.org/tools/impermanent-loss-calculator/}
\end{itemize}

\begin{rulebox}[Watch instead]
\begin{itemize}
\item ``What is an Automated Market Maker? (Liquidity Pool Algorithm)'' --- Whiteboard Crypto
(10:25). The CPMM story version of \S\ref{sec:cpmm}.
\url{https://www.youtube.com/watch?v=1PbZMudPP5E} (accessed 2026-10-11)
\item ``What is Impermanent Loss in Crypto? (Animated + Examples)'' --- Whiteboard Crypto
(8:17). Two worked IL examples; matches \S\ref{sec:il}.
\url{https://www.youtube.com/watch?v=_m6Mowq3Ptk} (accessed 2026-10-11)
\end{itemize}
\end{rulebox}
